\chapter{N DIMENSIONAL SPACES.}

We have already had the opportunity of saying how the development of the analytic Geometry of the plane and space led mathematicians to introduce the notion of n dimensional space, which supplied them with a geometric language which was extremely convenient for expressing in a concise and simple way theorems of algebra concerning equations with an arbitrary number of variables, and notably all the general results of linear Algebra (see p.~61). But if, towards the middle of the XIXth century, this language had become current among numerous geometers, it remained purely conventional, and the absence of an ``intuitive'' representation of spaces of more than three dimensions seemed to forbid in these latter arguments ``by continuity'' that were allowed in the plane or in space based exclusively on ``intuition''. It is Riemann who was the first, in his research on Analysis situs and on the foundations of Geometry, who was bold to argue in that way by analogy with the case of three dimensional space (see p.~176); \(^{1}\) following him, numerous mathematicians set out to use, with great success, arguments of this nature, notably in the theory of algebraic functions of several complex variables. But the control of spatial intuition being very limited at that time, one could remain with good reason very sceptical as to the worth in proofs of such considerations, and only allow them purely heuristically, in as much as they made very plausible the exactness of certain theorems. It is thus that H. Poincaré, in his memoir of 1887 on the residues of double integrals of functions of two complex variables, avoids as far as possible all recourse to intuition in four dimensional space: ``As this hypergeometric language is repugnant still to many good minds'', he says, ``I will only make infrequent use of it''; the ``artifices'' that he employs to this end allow him to get back to topological arguments in three dimensional space, where he then hesitates no longer to appeal to intuition ({[}251 a{]}, v. III, pp.~443 ff.).

Elsewhere, the discoveries of Cantor, and notably the famous theorem establishing that R and \(R^{n}\) are equipotent (which seemed to put into question the very notion of dimension) \(^{2}\) showed that it was indispensable, in order to provide a solid basis for the arguments of Geometry and Topology, to free them entirely from all recourse to intuition. We have already said (cf.~p.~139) that this need is at the origin of the modern conception of general Topology; but even before the creation of this latter, a rigorous study had begun of the topology of number spaces and of their most immediate generalisations (the ``n dimensional varieties'') by methods arising mainly from the branch of Topology called ``combinatorial Topology'' or better ``algebraic Topology'', of which we can not speak here.
% semantic-fragments: legacy-input-seam
\relax{}
